\chapter{Metodologia}

\subsection{Arquitectura general}

El pipeline consta de cuatro componentes conectados en serie. La
figura~\ref{fig:pipeline} muestra el flujo desde la adquisicion hasta
la alerta.

\begin{figure}[ht]
\centering
\begin{tikzpicture}[
  node distance=12mm and 8mm,
  every node/.style={font=\small},
  block/.style={
    rectangle, rounded corners=4pt,
    draw=blue!50!black, fill=blue!8,
    minimum width=28mm, minimum height=10mm,
    align=center, font=\small\bfseries
  },
  acc/.style={
    trapezium, trapezium left angle=70, trapezium right angle=110,
    draw=orange!70!black, fill=orange!10,
    minimum width=20mm, align=center, font=\small
  },
  out/.style={
    rectangle, draw=green!50!black, fill=green!10,
    minimum width=22mm, minimum height=10mm, align=center
  },
  arrow/.style={-{Stealth[length=2.5mm]}, thick}
]
\node[block] (tw) {Gemelo\\ digital ligero};
\node[acc, right=of tw] (resid) {Residuos\\ $r_t = x_t - \hat{x}_t$};
\node[block, right=of resid] (ae) {Autoencoder\\ LSTM};
\node[acc, right=of ae] (score) {$e_w = \text{MSE}$};
\node[block, right=of score] (cp) {Conformal\\ prediction};
\node[out, right=of cp] (alert) {Alerta\\(FPR $\le \alpha$)};

\node[block, above=of tw] (scada) {SCADA\\1 Hz};
\node[block, above=of ae] (fed) {FedAvg\\3 SEs};

\draw[arrow] (scada) -- (tw);
\draw[arrow] (scada) -- (resid);
\draw[arrow] (tw) -- (resid);
\draw[arrow] (resid) -- (ae);
\draw[arrow] (fed) -- (ae);
\draw[arrow] (ae) -- (score);
\draw[arrow] (score) -- (cp);
\draw[arrow] (cp) -- (alert);
\end{tikzpicture}
\caption{Diagrama del pipeline de deteccion.}
\label{fig:pipeline}
\end{figure}

\subsection{Gemelo digital ligero}

\subsubsection{Estado estacionario}

En regimen permanente, las corrientes por linea son

\[
I_k = \frac{S_k}{\sqrt{3}\, V_{BT}\, V_{k,p}}, \qquad k = 1,\dots,4
\]

donde $S_k = \sqrt{P_k^2 + Q_k^2}$ es la potencia aparente de la
linea $k$ y $V_{k,p}$ es la tension en el nodo $p$ aguas abajo del
alimentador. La tension en $p$ se calcula por la caida en la linea:

\[
V_{k,p} = V_{BT} - \frac{R_{linea,k}\, P_k + X_{linea,k}\, Q_k}{S_{nom}}
\]

La potencia total demandada al lado AT es

\[
P_{AT} = \sum_k P_k + 2 P_{Cu,nom} \left(\frac{S_{tot}}{2 S_{nom}}\right)^2 + 2 P_{Fe}
\]

\subsubsection{Dinamica termica}

La temperatura del devanado sigue una dinamica de primer orden

\[
C_{term}\, \frac{dT_{dev}}{dt} = P_{Cu}(t) - h (T_{dev} - T_{amb})
\]

con $C_{term} = h \cdot \tau_{term}$ y $P_{Cu}(t) = P_{Cu,nom} (S_{pu})^2$
donde $S_{pu} = S_{tot} / (2 S_{nom})$.

\subsubsection{Dinamica del regulador de tension}

El regulador automatico de tension ajusta el tap con dinamica de
constante 5~s:

\[
\frac{dV_{BT}}{dt} = \frac{V^{ref} - V_{BT}}{5\,\text{s}}
\]

El modelo se implementa en Python (~150 lineas,~\ref{appx:codigo})
y corre a 1~Hz sin problemas en CPU.

\subsection{Calculo de residuos}

El residuo multivariado en el tiempo $t$ es

\[
r_t = x_t^{obs} - \hat{x}_t^{gemelo} \in \mathbb{R}^{22}
\]

Bajo operacion normal, $r_t \approx \mathcal{N}(0, \Sigma)$ donde
$\Sigma$ es la covarianza empirica del ruido de medicion. Bajo
falla, $r_t$ exhibe una firma caracteristica:

\begin{itemize}[leftmargin=*,itemsep=2pt]
  \item F1 (cortocircuito): pico de corriente L1, depresion de
        tension AT/BT1, aumento de THD.
  \item F2 (sobrecarga): corriente L2 elevada con rampa sostenida.
  \item F3 (sobretension AT): tension AT por encima de 1.07~pu,
        frecuencia elevada.
  \item F4 (degradacion): deriva lenta del THD.
  \item F5 (PMU): frecuencia congelada y ROCOF nulo.
  \item F6 (desbalance): corrientes L1 y L3 opuestas.
  \item F7 (buje): pulsos esporadicos en THD.
  \item F8 (clima): tension AT cae, corriente sube, frecuencia cae.
\end{itemize}

El residuo es la materia prima del detector no supervisado. Es
\emph{robusto al punto de operacion}: si la carga sube, el gemelo
sube la corriente esperada y el residuo sigue siendo pequeno.

\subsection{Autoencoder LSTM}

El detector es un autoencoder LSTM con la arquitectura de la
figura~\ref{fig:ae}.

\begin{figure}[ht]
\centering
\begin{tikzpicture}[
  every node/.style={font=\small},
  box/.style={
    rectangle, draw=blue!60!black, fill=blue!10,
    minimum width=22mm, minimum height=8mm, align=center
  },
  arrow/.style={-{Stealth[length=2.5mm]}, thick}
]
\node[box] (in) {Input\\$30 \times 22$};
\node[box, right=4mm of in] (e1) {LSTM 64};
\node[box, right=2mm of e1] (e2) {LSTM 32};
\node[box, right=2mm of e2] (z) {Latent 16};
\node[box, right=4mm of z] (d1) {Repeat 30};
\node[box, right=2mm of d1] (d2) {LSTM 32};
\node[box, right=2mm of d2] (d3) {LSTM 64};
\node[box, right=2mm of d3] (out) {Dense 22};

\draw[arrow] (in) -- (e1);
\draw[arrow] (e1) -- (e2);
\draw[arrow] (e2) -- (z);
\draw[arrow] (z) -- (d1);
\draw[arrow] (d1) -- (d2);
\draw[arrow] (d2) -- (d3);
\draw[arrow] (d3) -- (out);

\draw[arrow, dashed, gray] ([yshift=-3mm]in.south) -- node[below,font=\tiny,text=gray]{MSE loss} ([yshift=-3mm]out.south);
\end{tikzpicture}
\caption{Arquitectura del autoencoder LSTM.}
\label{fig:ae}
\end{figure}

La entrada es una ventana $w$ de $T=30$ muestras de los 22 features
normalizados ($\mu, \sigma$ calculados sobre el conjunto de
entrenamiento normal). El encoder produce un vector latente de 16
dimensiones, que el decoder repite $T$ veces para reconstruir la
ventana completa. La loss es MSE.

El modelo tiene $\approx 53{,}000$ parametros. Se entrena con
optimizer Adam ($lr=10^{-3}$, weight decay $10^{-5}$),
CosineAnnealingLR y early stopping sobre validacion (paciencia 5).
30 epocas bastan para converger.

\subsection{Conformal prediction}

El umbral conformal se calcula sobre el conjunto de calibracion
(50\% de las ventanas de train) con $\alpha = 0.05$:

\[
\hat{q} = \text{Quantil}_{n_{cal}+1}^{\lceil 0.95 \cdot (n_{cal}+1) \rceil}\, (e_w^{cal})
\]

Toda ventana de test con reconstruction error $e_w^{test} > \hat{q}$
genera una alerta. Bajo intercambiabilidad,

\[
\mathbb{P}\{e_w^{test} > \hat{q}\} \le 0.05
\]

lo que da una \emph{tasa de falsa alarma garantizada} del 5\%.

\subsubsection{Por que calibrar sobre el train y no sobre un conjunto separado}

Por restricciones de tamano muestral, las primeras 50\% de las
ventanas de train se usan como calibracion. Para uso en produccion
se recomienda reservar un conjunto separado de $\sim 100{,}000$
muestras normales. La validacion del supuesto de intercambiabilidad
se hace graficando la CDF empirica del reconstruction error sobre el
test (capitulo 5).

\subsection{Aprendizaje federado}

Se construyen tres subestaciones virtuales (configuracion reducida para
este prototipo; en produccion se usan cinco o mas), cada una con un perfil de
carga y catalogo de fallas ligeramente distinto (variabilidad
geografica). El entrenamiento sigue FedAvg:

\begin{enumerate}[leftmargin=*,itemsep=2pt]
  \item Inicializar $\theta_0$ en el servidor.
  \item Para cada ronda $t = 1, \dots, T$:
  \begin{enumerate}
    \item Enviar $\theta_t$ a cada cliente $k$.
    \item Entrenar $\theta_t$ sobre $D_k$ durante $E$ epocas
          locales, obtener $\theta^k_{t+1}$.
    \item Agregar: $\theta_{t+1} = \sum_k \frac{n_k}{N} \theta^k_{t+1}$.
  \end{enumerate}
\end{enumerate}

Parametros: $T = 10$ rondas, $E = 3$ epocas locales, $lr = 10^{-3}$.
No se usa ninguna tecnica de privacidad diferencial (DP) en esta
version por simplicidad, pero se deja el hook para agregarla en
produccion.

\subsection{Priorizacion post-evento}

Tras un frente climatico detectado, el sistema emite un orden de
despacho de cuadrillas por subestacion. El score de prioridad es:

\[
P_i = w_1 \cdot \mathbb{1}\{alerta_i > \hat{q}\}
    + w_2 \cdot \mathbb{1}\{T_{dev,i} > T_{crit}\}
    + w_3 \cdot \mathbb{1}\{f_i < f_{crit}\}
    + w_4 \cdot \text{criticidad}_i
\]

donde los $w_j$ son hiperparametros que la distribuidora fija. Una
heuristica razonable es $w = (0.4, 0.2, 0.2, 0.2)$. Las SE con mayor
$P_i$ son atendidas primero.

\subsection{Evaluacion}

\subsubsection{Metricas}

\begin{itemize}[leftmargin=*,itemsep=2pt]
  \item \textbf{FPR} (False Positive Rate): proporcion de ventanas
        normales flaggeadas como falla.
  \item \textbf{TPR / Recall}: proporcion de ventanas con falla
        detectadas.
  \item \textbf{Precision}: proporcion de alertas que son falla real.
  \item \textbf{F1}: media armonica de precision y recall.
  \item \textbf{AUC}: area bajo la curva ROC, umbral-agnostico.
  \item \textbf{Latencia de deteccion}: tiempo entre el inicio de la
        falla y la primera alerta.
\end{itemize}

\subsubsection{Comparacion de detectores}

Se comparan tres detectores:
\begin{enumerate}[leftmargin=*,itemsep=2pt]
  \item \emph{Estatico 3-sigma}: umbral fijo sobre $|r_{V_{AT}}|$.
        \emph{Baseline} que muestra la necesidad de modelos mas
        sofisticados.
  \item \emph{AE-LSTM centralizado}: autoencoder entrenado sobre
        todos los datos centralizados. \emph{Techo} teorico de
        rendimiento.
  \item \emph{AE-LSTM federado (3 SEs)}: el modelo propuesto,
        preservando privacidad.
\end{enumerate}

\subsubsection{Cumplimiento de la garantia conformal}

El conformal prediction entrega una cota \emph{teorica} sobre el FPR.
La verificacion empirica se hace midiendo el FPR sobre el conjunto
de test: si el FPR empirico es $\le \alpha$, la garantia se cumple
en la muestra. Los resultados del capitulo siguiente muestran que
el FPR empirico es de 5.4\%, ligeramente por encima del objetivo del
5\%. Esta pequena brecha es esperada y se discute en el capitulo 6.